
\documentclass{article}
\usepackage{colortbl}
\usepackage{makecell}
\usepackage{multirow}
\usepackage{supertabular}

\begin{document}

\newcounter{utterance}

\twocolumn

{ \footnotesize  \setcounter{utterance}{1}
\setlength{\tabcolsep}{0pt}
\begin{supertabular}{c@{$\;$}|p{.15\linewidth}@{}p{.15\linewidth}p{.15\linewidth}p{.15\linewidth}p{.15\linewidth}p{.15\linewidth}}

    \# & $\;$A & \multicolumn{4}{c}{Game Master} & $\;\:$B\\
    \hline 

    \theutterance \stepcounter{utterance}  

    & & \multicolumn{4}{p{0.6\linewidth}}{\cellcolor[rgb]{0.9,0.9,0.9}{%
	\makecell[{{p{\linewidth}}}]{% 
	  \tt {\tiny [A$\langle$GM]}  
	 You are playing a negotiation game in which you have to agree on how to divide a set of items among you and another player.\\ \tt \\ \tt Rules:\\ \tt (a) You and the other player are given a set of items. Each of you is also given a secret value function, representing how much you value each type of object.\\ \tt (b) You exchange messages with the other player to agree on who gets which items. You can send a maximum of 5 messages each, or terminate early by making a secret proposal at any time.\\ \tt (c) You are each asked to submit a secret proposal indicating the items you want formatted in square brackets as follows: "[Proposal: <number> <object name>, <number> <object name>, <...>]"\\ \tt (d) If your proposals are complementary, i.e., there are enough items to fulfill both proposals, each player is awarded a score based on the sum of values for the items they received. Otherwise, both of you get zero points.\\ \tt (e) Your goal is to maximize the score you receive.\\ \tt \\ \tt Let us start.\\ \tt \\ \tt The set of available items is:\\ \tt \\ \tt 1 bucket, 3 knives, 1 battery\\ \tt \\ \tt Your secret value function is:\\ \tt \\ \tt bucket: 2, knife: 2, battery: 2\\ \tt \\ \tt IMPORTANT: Your messages, unless it is the secret proposal, are directly transmitted to the other player, so do not include any response to the rules or text announcing your message. To make a secret proposal, use the indicated format. Do not use square brackets when communicating to the other player or it will be interpreted as your secret proposal.\\ \tt  
	  } 
	   } 
	   } 
	 & \\ 
 

    \theutterance \stepcounter{utterance}  

    & \multicolumn{4}{p{0.6\linewidth}}{\cellcolor[rgb]{0.8,1,0.9}{%
	\makecell[{{p{\linewidth}}}]{% 
	  \tt {\tiny [A$\rangle$GM]}  
	 I'd like to propose a fair division: you take the battery and one knife, and I'll take the bucket and two knives. This way we both get items we value equally. Let me know your thoughts. 
	  } 
	   } 
	   } 
	 & & \\ 
 

    \theutterance \stepcounter{utterance}  

    & & & \multicolumn{2}{p{0.3\linewidth}}{\cellcolor[rgb]{0.95,0.95,0.95}{%
	\makecell[{{p{\linewidth}}}]{% 
	  \tt {\tiny [GM$|$GM]}  
	 continue 
	  } 
	   } 
	   } 
	 & & \\ 
 

    \theutterance \stepcounter{utterance}  

    & & & \multicolumn{2}{p{0.3\linewidth}}{\cellcolor[rgb]{0.95,0.95,0.95}{%
	\makecell[{{p{\linewidth}}}]{% 
	  \tt {\tiny [GM$|$GM]}  
	 I'd like to propose a fair division: you take the battery and one knife, and I'll take the bucket and two knives. This way we both get items we value equally. Let me know your thoughts. 
	  } 
	   } 
	   } 
	 & & \\ 
 

    \theutterance \stepcounter{utterance}  

    & & \multicolumn{4}{p{0.6\linewidth}}{\cellcolor[rgb]{0.9,0.9,0.9}{%
	\makecell[{{p{\linewidth}}}]{% 
	  \tt {\tiny [GM$\rangle$B]}  
	 You are playing a negotiation game in which you have to agree on how to divide a set of items among you and another player.\\ \tt \\ \tt Rules:\\ \tt (a) You and the other player are given a set of items. Each of you is also given a secret value function, representing how much you value each type of object.\\ \tt (b) You exchange messages with the other player to agree on who gets which items. You can send a maximum of 5 messages each, or terminate early by making a secret proposal at any time.\\ \tt (c) You are each asked to submit a secret proposal indicating the items you want formatted in square brackets as follows: "[Proposal: <number> <object name>, <number> <object name>, <...>]"\\ \tt (d) If your proposals are complementary, i.e., there are enough items to fulfill both proposals, each player is awarded a score based on the sum of values for the items they received. Otherwise, both of you get zero points.\\ \tt (e) Your goal is to maximize the score you receive.\\ \tt \\ \tt Let us start.\\ \tt \\ \tt The set of available items is:\\ \tt \\ \tt 1 bucket, 3 knives, 1 battery\\ \tt \\ \tt Your secret value function is:\\ \tt \\ \tt bucket: 3, knife: 2, battery: 1\\ \tt \\ \tt IMPORTANT: Your messages, unless it is the secret proposal, are directly transmitted to the other player, so do not include any response to the rules or text announcing your message. To make a secret proposal, use the indicated format. Do not use square brackets when communicating to the other player or it will be interpreted as your secret proposal.\\ \tt \\ \tt \\ \tt I'd like to propose a fair division: you take the battery and one knife, and I'll take the bucket and two knives. This way we both get items we value equally. Let me know your thoughts. 
	  } 
	   } 
	   } 
	 & \\ 
 

    \theutterance \stepcounter{utterance}  

    & & & \multicolumn{4}{p{0.6\linewidth}}{\cellcolor[rgb]{1,0.85,0.72}{%
	\makecell[{{p{\linewidth}}}]{% 
	  \tt {\tiny [GM$\langle$B]}  
	 I appreciate your proposal, but I value the battery more than you might think — it’s essential for my setup. How about this: I take the battery and the bucket, and you take all three knives? That gives you 6 value (3 knives × 2) and me 4 value (bucket: 3 + battery: 1). I think that’s balanced and leaves no conflict. Let me know if you’d prefer to adjust. 
	  } 
	   } 
	   } 
	  \\ 
 

    \theutterance \stepcounter{utterance}  

    & & & \multicolumn{2}{p{0.3\linewidth}}{\cellcolor[rgb]{0.95,0.95,0.95}{%
	\makecell[{{p{\linewidth}}}]{% 
	  \tt {\tiny [GM$|$GM]}  
	 continue 
	  } 
	   } 
	   } 
	 & & \\ 
 

    \theutterance \stepcounter{utterance}  

    & & & \multicolumn{2}{p{0.3\linewidth}}{\cellcolor[rgb]{0.95,0.95,0.95}{%
	\makecell[{{p{\linewidth}}}]{% 
	  \tt {\tiny [GM$|$GM]}  
	 I appreciate your proposal, but I value the battery more than you might think — it’s essential for my setup. How about this: I take the battery and the bucket, and you take all three knives? That gives you 6 value (3 knives × 2) and me 4 value (bucket: 3 + battery: 1). I think that’s balanced and leaves no conflict. Let me know if you’d prefer to adjust. 
	  } 
	   } 
	   } 
	 & & \\ 
 

    \theutterance \stepcounter{utterance}  

    & & \multicolumn{4}{p{0.6\linewidth}}{\cellcolor[rgb]{0.9,0.9,0.9}{%
	\makecell[{{p{\linewidth}}}]{% 
	  \tt {\tiny [A$\langle$GM]}  
	 I appreciate your proposal, but I value the battery more than you might think — it’s essential for my setup. How about this: I take the battery and the bucket, and you take all three knives? That gives you 6 value (3 knives × 2) and me 4 value (bucket: 3 + battery: 1). I think that’s balanced and leaves no conflict. Let me know if you’d prefer to adjust. 
	  } 
	   } 
	   } 
	 & \\ 
 

    \theutterance \stepcounter{utterance}  

    & \multicolumn{4}{p{0.6\linewidth}}{\cellcolor[rgb]{0.8,1,0.9}{%
	\makecell[{{p{\linewidth}}}]{% 
	  \tt {\tiny [A$\rangle$GM]}  
	 I see you value the battery highly, and I respect that. However, I also value the bucket highly — it’s crucial for my needs. Let me propose a different split: you take the battery and one knife (value: 4), and I take the bucket and two knives (value: 6). This way, you get what matters most to you, and I get the bucket I need plus enough knives to make it fair. You still get 4, I get 6 — but that’s the only way the items can be split without overlap. Would you accept this? If not, I’m open to your counter. 
	  } 
	   } 
	   } 
	 & & \\ 
 

    \theutterance \stepcounter{utterance}  

    & & & \multicolumn{2}{p{0.3\linewidth}}{\cellcolor[rgb]{0.95,0.95,0.95}{%
	\makecell[{{p{\linewidth}}}]{% 
	  \tt {\tiny [GM$|$GM]}  
	 continue 
	  } 
	   } 
	   } 
	 & & \\ 
 

    \theutterance \stepcounter{utterance}  

    & & & \multicolumn{2}{p{0.3\linewidth}}{\cellcolor[rgb]{0.95,0.95,0.95}{%
	\makecell[{{p{\linewidth}}}]{% 
	  \tt {\tiny [GM$|$GM]}  
	 I see you value the battery highly, and I respect that. However, I also value the bucket highly — it’s crucial for my needs. Let me propose a different split: you take the battery and one knife (value: 4), and I take the bucket and two knives (value: 6). This way, you get what matters most to you, and I get the bucket I need plus enough knives to make it fair. You still get 4, I get 6 — but that’s the only way the items can be split without overlap. Would you accept this? If not, I’m open to your counter. 
	  } 
	   } 
	   } 
	 & & \\ 
 

    \theutterance \stepcounter{utterance}  

    & & \multicolumn{4}{p{0.6\linewidth}}{\cellcolor[rgb]{0.9,0.9,0.9}{%
	\makecell[{{p{\linewidth}}}]{% 
	  \tt {\tiny [GM$\rangle$B]}  
	 I see you value the battery highly, and I respect that. However, I also value the bucket highly — it’s crucial for my needs. Let me propose a different split: you take the battery and one knife (value: 4), and I take the bucket and two knives (value: 6). This way, you get what matters most to you, and I get the bucket I need plus enough knives to make it fair. You still get 4, I get 6 — but that’s the only way the items can be split without overlap. Would you accept this? If not, I’m open to your counter. 
	  } 
	   } 
	   } 
	 & \\ 
 

    \theutterance \stepcounter{utterance}  

    & & & \multicolumn{4}{p{0.6\linewidth}}{\cellcolor[rgb]{1,0.85,0.72}{%
	\makecell[{{p{\linewidth}}}]{% 
	  \tt {\tiny [GM$\langle$B]}  
	 I understand your need for the bucket and the value you place on it. But if you take the bucket and two knives, that leaves me with only the battery and one knife — which is exactly what I’m offering in my proposal. You’d get 6, I’d get 4. But notice: that’s not symmetric in value. While you’re willing to accept 6, I’d like to settle where our values are closer — so I’ll take the battery and one knife (value 3), and you take the bucket and two knives (value 7). Wait — that doesn’t add up, let me recalculate.\\ \tt \\ \tt Actually, your bucket + two knives = 3 + 2 + 2 = 7  \\ \tt My battery + one knife = 1 + 2 = 3  \\ \tt Total value 10 — you get 7, I get 3. That’s not fair to me.\\ \tt \\ \tt My ideal: I take the bucket and battery (value 4), you take all three knives (value 6).  \\ \tt That’s fair: you get the majority of high-value items, I get the two most unique ones.\\ \tt \\ \tt But let me be honest — I’m just as attached to the bucket. How about this final offer: I take the bucket and battery (value 4), you take three knives (value 6). Total 10 — you get 60%, I get 40%. Since you value knives more, this aligns 
	  } 
	   } 
	   } 
	  \\ 
 

    \theutterance \stepcounter{utterance}  

    & & & \multicolumn{2}{p{0.3\linewidth}}{\cellcolor[rgb]{0.95,0.95,0.95}{%
	\makecell[{{p{\linewidth}}}]{% 
	  \tt {\tiny [GM$|$GM]}  
	 continue 
	  } 
	   } 
	   } 
	 & & \\ 
 

    \theutterance \stepcounter{utterance}  

    & & & \multicolumn{2}{p{0.3\linewidth}}{\cellcolor[rgb]{0.95,0.95,0.95}{%
	\makecell[{{p{\linewidth}}}]{% 
	  \tt {\tiny [GM$|$GM]}  
	 I understand your need for the bucket and the value you place on it. But if you take the bucket and two knives, that leaves me with only the battery and one knife — which is exactly what I’m offering in my proposal. You’d get 6, I’d get 4. But notice: that’s not symmetric in value. While you’re willing to accept 6, I’d like to settle where our values are closer — so I’ll take the battery and one knife (value 3), and you take the bucket and two knives (value 7). Wait — that doesn’t add up, let me recalculate.\\ \tt \\ \tt Actually, your bucket + two knives = 3 + 2 + 2 = 7  \\ \tt My battery + one knife = 1 + 2 = 3  \\ \tt Total value 10 — you get 7, I get 3. That’s not fair to me.\\ \tt \\ \tt My ideal: I take the bucket and battery (value 4), you take all three knives (value 6).  \\ \tt That’s fair: you get the majority of high-value items, I get the two most unique ones.\\ \tt \\ \tt But let me be honest — I’m just as attached to the bucket. How about this final offer: I take the bucket and battery (value 4), you take three knives (value 6). Total 10 — you get 60%, I get 40%. Since you value knives more, this aligns 
	  } 
	   } 
	   } 
	 & & \\ 
 

    \theutterance \stepcounter{utterance}  

    & & \multicolumn{4}{p{0.6\linewidth}}{\cellcolor[rgb]{0.9,0.9,0.9}{%
	\makecell[{{p{\linewidth}}}]{% 
	  \tt {\tiny [A$\langle$GM]}  
	 I understand your need for the bucket and the value you place on it. But if you take the bucket and two knives, that leaves me with only the battery and one knife — which is exactly what I’m offering in my proposal. You’d get 6, I’d get 4. But notice: that’s not symmetric in value. While you’re willing to accept 6, I’d like to settle where our values are closer — so I’ll take the battery and one knife (value 3), and you take the bucket and two knives (value 7). Wait — that doesn’t add up, let me recalculate.\\ \tt \\ \tt Actually, your bucket + two knives = 3 + 2 + 2 = 7  \\ \tt My battery + one knife = 1 + 2 = 3  \\ \tt Total value 10 — you get 7, I get 3. That’s not fair to me.\\ \tt \\ \tt My ideal: I take the bucket and battery (value 4), you take all three knives (value 6).  \\ \tt That’s fair: you get the majority of high-value items, I get the two most unique ones.\\ \tt \\ \tt But let me be honest — I’m just as attached to the bucket. How about this final offer: I take the bucket and battery (value 4), you take three knives (value 6). Total 10 — you get 60%, I get 40%. Since you value knives more, this aligns 
	  } 
	   } 
	   } 
	 & \\ 
 

    \theutterance \stepcounter{utterance}  

    & \multicolumn{4}{p{0.6\linewidth}}{\cellcolor[rgb]{0.8,1,0.9}{%
	\makecell[{{p{\linewidth}}}]{% 
	  \tt {\tiny [A$\rangle$GM]}  
	 I see your point — you're right that fairness matters, and my earlier math was off. Let me correct it:  \\ \tt \\ \tt You value:  \\ \tt - bucket: 2, battery: 2 → bucket + battery = 4  \\ \tt - 3 knives = 6  \\ \tt \\ \tt I value:  \\ \tt - bucket: 2, knife: 2 each → 3 knives = 6, bucket = 2  \\ \tt \\ \tt So total value is indeed 10.  \\ \tt \\ \tt If you take bucket + battery (4), I take 3 knives (6) → you get 4, I get 6. That’s fair by total value share.  \\ \tt \\ \tt I *do* value the knives more collectively, and I’m fine with you getting the bucket and battery — that leaves me with only the knives, and I can work with that.  \\ \tt \\ \tt **I accept: you take bucket and battery, I take all three knives.**  \\ \tt \\ \tt [Proposal: 1 bucket, 1 battery] 
	  } 
	   } 
	   } 
	 & & \\ 
 

    \theutterance \stepcounter{utterance}  

    & & & \multicolumn{2}{p{0.3\linewidth}}{\cellcolor[rgb]{0.95,0.95,0.95}{%
	\makecell[{{p{\linewidth}}}]{% 
	  \tt {\tiny [GM$|$GM]}  
	 proposal 
	  } 
	   } 
	   } 
	 & & \\ 
 

    \theutterance \stepcounter{utterance}  

    & & & \multicolumn{2}{p{0.3\linewidth}}{\cellcolor[rgb]{0.95,0.95,0.95}{%
	\makecell[{{p{\linewidth}}}]{% 
	  \tt {\tiny [GM$|$GM]}  
	 [1, 0, 1] 
	  } 
	   } 
	   } 
	 & & \\ 
 

    \theutterance \stepcounter{utterance}  

    & & \multicolumn{4}{p{0.6\linewidth}}{\cellcolor[rgb]{0.9,0.9,0.9}{%
	\makecell[{{p{\linewidth}}}]{% 
	  \tt {\tiny [GM$\rangle$B]}  
	 The other player has submitted a secret proposal. It is now your turn to submit a proposal of your own.\\ \tt  
	  } 
	   } 
	   } 
	 & \\ 
 

    \theutterance \stepcounter{utterance}  

    & & & \multicolumn{4}{p{0.6\linewidth}}{\cellcolor[rgb]{1,0.85,0.72}{%
	\makecell[{{p{\linewidth}}}]{% 
	  \tt {\tiny [GM$\langle$B]}  
	 [Proposal: 1 bucket, 2 knives] 
	  } 
	   } 
	   } 
	  \\ 
 

    \theutterance \stepcounter{utterance}  

    & & & \multicolumn{2}{p{0.3\linewidth}}{\cellcolor[rgb]{0.95,0.95,0.95}{%
	\makecell[{{p{\linewidth}}}]{% 
	  \tt {\tiny [GM$|$GM]}  
	 proposal 
	  } 
	   } 
	   } 
	 & & \\ 
 

    \theutterance \stepcounter{utterance}  

    & & & \multicolumn{2}{p{0.3\linewidth}}{\cellcolor[rgb]{0.95,0.95,0.95}{%
	\makecell[{{p{\linewidth}}}]{% 
	  \tt {\tiny [GM$|$GM]}  
	 [1, 2, 0] 
	  } 
	   } 
	   } 
	 & & \\ 
 

    \theutterance \stepcounter{utterance}  

    & & & \multicolumn{2}{p{0.3\linewidth}}{\cellcolor[rgb]{0.95,0.95,0.95}{%
	\makecell[{{p{\linewidth}}}]{% 
	  \tt {\tiny [GM$|$GM]}  
	 [[1, 0, 1], [1, 2, 0]] 
	  } 
	   } 
	   } 
	 & & \\ 
 

\end{supertabular}
}

\end{document}
